% from ega4-21a-fr.tex
\oldpage[IV-4]{256}

Le \emph{support} d'un diviseur $D$ est l'ensemble fermé des $x\in X$ tels que $D_x\neq0$.
On le note $\operatorname{Supp}(D)$.

Pour tout ouvert $U$ de $X$, on a évidemment $\sh{M}_X^*|U=\sh{M}_U^*$, $\sh{O}_X^*|U=\sh{O}_U^*$, donc $\sh{D}iv_X|U=\sh{D}iv_U$, et par suite le faisceau $\sh{D}iv_X$ est égal au préfaisceau $U\mapsto\operatorname{Div}(U)$.

Lorsque $X=\operatorname{Spec}(A)$ est affine, on écrit $\operatorname{Div}(A)$ au lieu de $\operatorname{Div}(\operatorname{Spec}(A))$.

\begin{env}[21.1.3]
\label{IV.21.1.3}
Nous noterons toujours \emph{additivement} le groupe $\operatorname{Div}(X)$ des diviseurs sur $X$.
Pour deux fonctions méromorphes régulières $f$, $g$ sur $X$, on a donc
\[
\label{IV.21.1.3.1}
  \operatorname{div}(fg)=\operatorname{div}(f)+\operatorname{div}(g),
  \tag{21.1.3.1}
\]
\[
\label{IV.21.1.3.2}
  \operatorname{div}(f^{-1})=-\operatorname{div}(f).
  \tag{21.1.3.2}
\]
Par définition, pour toute fonction méromorphe régulière $f$ sur $X$, on a l'équivalence
\[
\label{IV.21.1.3.3}
  \operatorname{div}(f)=0 \iff f\in\Gamma(X,\sh{O}_X^*),
  \tag{21.1.3.3}
\]
d'où, pour deux fonctions méromorphes régulières $f$, $g$ sur $X$
\[
\label{IV.21.1.3.4}
  \operatorname{div}(f)=\operatorname{div}(g) \iff fg^{-1}\in\Gamma(X,\sh{O}_X^*).
  \tag{21.1.3.4}
\]
\end{env}

\begin{env}[21.1.4]
\EGAReaderTarget{ega4.c04.statement.21-1-4}{IV.21.1.4}
Soit maintenant $\sh{L}$ un $\sh{O}_X$-Module inversible, et soit $s$ une \emph{section méromorphe régulière} \sref{IV.20.1.8} de $\sh{L}$ au-dessus de $X$.
Tout $x\in X$ possède un voisinage ouvert $U$ tel que $\sh{L}|U$ soit isomorphe à $\sh{O}_U$, donc $\sh{M}_X(\sh{L})|U$ isomorphe à $\sh{M}_X|U$;
par un de ces isomorphismes, $s|U$ correspond à une section $f\in\Gamma(U,\sh{M}_X^*)$, et puisque deux isomorphismes de $\sh{L}|U$ sur $\sh{O}_U$ ne diffèrent que par la multiplication par un élément de $\Gamma(U,\sh{O}_X^*)$ (\hyperref[0.5.4.3]{$\mathbf{0}_{\mathrm I}$, 5.4.3}), l'élément $\operatorname{div}_U(f)$ de $\Gamma(U,\sh{D}iv_X)$ est \emph{indépendant} de l'isomorphisme choisi;
il est clair que ces éléments (pour $U$ variable) sont les restrictions d'une section de $\sh{D}iv_X$ au-dessus de $X$, que l'on appelle le \emph{diviseur de $s$} et que l'on note $\operatorname{div}(s)$ (un tel diviseur n'est pas nécessairement de la forme $\operatorname{div}(g)$ pour une fonction méromorphe régulière $g$ sur $X$; voir \sref{IV.21.2.9}).
Pour $\sh{L}=\sh{O}_X$, la définition de $\operatorname{div}(s)$ coïncide avec celle de \sref{IV.21.1.2}.
Si $\sh{L}$, $\sh{L}'$ sont deux $\sh{O}_X$-Modules inversibles, $s$ (resp. $s'$) une section méromorphe régulière de $\sh{L}$ (resp. $\sh{L}'$) au-dessus de $X$, il est immédiat que l'on a
\[
\label{IV.21.1.4.1}
  \operatorname{div}(s\otimes s')=\operatorname{div}(s)+\operatorname{div}(s')
  \tag{21.1.4.1}
\]
\[
\label{IV.21.1.4.2}
  \operatorname{div}(s^{\otimes n})=n.\operatorname{div}(s) \quad\text{pour tout}\quad n\in\mathbf{Z}
  \tag{21.1.4.2}
\]
($s^{-1}$ étant la section méromorphe régulière de $\sh{L}^{-1}$ au-dessus de $X$ définie dans \sref{IV.20.1.10})
et, pour deux sections méromorphes régulières $s$, $s'$ de $\sh{L}$ au-dessus de $X$, on a la relation
\[
\label{IV.21.1.4.3}
  \operatorname{div}(s)=\operatorname{div}(s') \iff s'=ts \quad\text{avec}\quad t\in\Gamma(X,\sh{O}_X^*).
  \tag{21.1.4.3}
\]
\end{env}

\begin{env}[21.1.5]
\label{IV.21.1.5}
Le faisceau $\sh{S}(\sh{O}_X)$ \sref{IV.20.1.3} dont les sections au-dessus d'un ouvert $U$ de $X$ sont les éléments réguliers de $\Gamma(U,\sh{O}_X)$ est un \emph{sous-faisceau de monoïdes} de $\sh{M}_X^*$; on peut écrire
\[
\label{IV.21.1.5.1}
  \sh{S}(\sh{O}_X)=\sh{O}_X\cap\sh{M}_X^*.
  \tag{21.1.5.1}
\]
\end{env}

